\documentclass[11pt]{article}
\usepackage{float}
% NOTE: Add in the relevant information to the commands below; or, if you'll be using the same information frequently, add these commands at the top of paolo-pset.tex file. 
\newcommand{\name}{Agustin Esteva}
\newcommand{\email}{aesteva@uchicago.edu}
\newcommand{\classnum}{13310}
\newcommand{\subject}{SSI: Formal Theory}
\newcommand{\instructors}{Zhaotian Luo}
\newcommand{\assignment}{Problem Set 1}
\newcommand{\semester}{Spring 2025}
\newcommand{\duedate}{23-04-25}
\newcommand{\bA}{\mathbf{A}}
\newcommand{\bB}{\mathbf{B}}
\newcommand{\bC}{\mathbf{C}}
\newcommand{\bD}{\mathbf{D}}
\newcommand{\bE}{\mathbf{E}}
\newcommand{\bF}{\mathbf{F}}
\newcommand{\bG}{\mathbf{G}}
\newcommand{\bH}{\mathbf{H}}
\newcommand{\bI}{\mathbf{I}}
\newcommand{\bJ}{\mathbf{J}}
\newcommand{\bK}{\mathbf{K}}
\newcommand{\bL}{\mathbf{L}}
\newcommand{\bM}{\mathbf{M}}
\newcommand{\bN}{\mathbf{N}}
\newcommand{\bO}{\mathbf{O}}
\newcommand{\bP}{\mathbf{P}}
\newcommand{\bQ}{\mathbf{Q}}
\newcommand{\bR}{\mathbf{R}}
\newcommand{\bS}{\mathbf{S}}
\newcommand{\bT}{\mathbf{T}}
\newcommand{\bU}{\mathbf{U}}
\newcommand{\bV}{\mathbf{V}}
\newcommand{\bW}{\mathbf{W}}
\newcommand{\bX}{\mathbf{X}}
\newcommand{\bY}{\mathbf{Y}}
\newcommand{\bZ}{\mathbf{Z}}
\newcommand{\Var}{\text{Var}}

%% blackboard bold math capitals
\newcommand{\bbA}{\mathbb{A}}
\newcommand{\bbB}{\mathbb{B}}
\newcommand{\bbC}{\mathbb{C}}
\newcommand{\bbD}{\mathbb{D}}
\newcommand{\bbE}{\mathbb{E}}
\newcommand{\bbF}{\mathbb{F}}
\newcommand{\bbG}{\mathbb{G}}
\newcommand{\bbH}{\mathbb{H}}
\newcommand{\bbI}{\mathbb{I}}
\newcommand{\bbJ}{\mathbb{J}}
\newcommand{\bbK}{\mathbb{K}}
\newcommand{\bbL}{\mathbb{L}}
\newcommand{\bbM}{\mathbb{M}}
\newcommand{\bbN}{\mathbb{N}}
\newcommand{\bbO}{\mathbb{O}}
\newcommand{\bbP}{\mathbb{P}}
\newcommand{\bbQ}{\mathbb{Q}}
\newcommand{\bbR}{\mathbb{R}}
\newcommand{\bbS}{\mathbb{S}}
\newcommand{\bbT}{\mathbb{T}}
\newcommand{\bbU}{\mathbb{U}}
\newcommand{\bbV}{\mathbb{V}}
\newcommand{\bbW}{\mathbb{W}}
\newcommand{\bbX}{\mathbb{X}}
\newcommand{\bbY}{\mathbb{Y}}
\newcommand{\bbZ}{\mathbb{Z}}

%% script math capitals
\newcommand{\sA}{\mathscr{A}}
\newcommand{\sB}{\mathscr{B}}
\newcommand{\sC}{\mathscr{C}}
\newcommand{\sD}{\mathscr{D}}
\newcommand{\sE}{\mathscr{E}}
\newcommand{\sF}{\mathscr{F}}
\newcommand{\sG}{\mathscr{G}}
\newcommand{\sH}{\mathscr{H}}
\newcommand{\sI}{\mathscr{I}}
\newcommand{\sJ}{\mathscr{J}}
\newcommand{\sK}{\mathscr{K}}
\newcommand{\sL}{\mathscr{L}}
\newcommand{\sM}{\mathscr{M}}
\newcommand{\sN}{\mathscr{N}}
\newcommand{\sO}{\mathscr{O}}
\newcommand{\sP}{\mathscr{P}}
\newcommand{\sQ}{\mathscr{Q}}
\newcommand{\sR}{\mathscr{R}}
\newcommand{\sS}{\mathscr{S}}
\newcommand{\sT}{\mathscr{T}}
\newcommand{\sU}{\mathscr{U}}
\newcommand{\sV}{\mathscr{V}}
\newcommand{\sW}{\mathscr{W}}
\newcommand{\sX}{\mathscr{X}}
\newcommand{\sY}{\mathscr{Y}}
\newcommand{\sZ}{\mathscr{Z}}


\renewcommand{\emptyset}{\O}

\newcommand{\abs}[1]{\lvert #1 \rvert}
\newcommand{\norm}[1]{\lVert #1 \rVert}
\newcommand{\sm}{\setminus}


\newcommand{\sarr}{\rightarrow}
\newcommand{\arr}{\longrightarrow}

% NOTE: Defining collaborators is optional; to not list collaborators, comment out the line below.
%\newcommand{\collaborators}{Alyssa P. Hacker (\texttt{aphacker}), Ben Bitdiddle (\texttt{bitdiddle})}

\input{paolo-pset.tex}

% NOTE: To compile a version of this pset without problems, solutions, or reflections, uncomment the relevant line below.

%\excludeversion{problem}
%\excludeversion{solution}
%\excludeversion{reflection}

\begin{document}	
	
	% Use the \psetheader command at the beginning of a pset. 
	\psetheader
\section*{Problem 1}

\textbf{Payoff Matrix:}

\[
\begin{array}{c|c|c}
1/2 & \text{Work} & \text{Shirk} \\
\hline
\text{Work} & r - c, \theta - c & -c, 0 \\
\text{Shirk} & 0, -c & 0, 0 \\
\end{array}
\]

Player 1's value of collaboration \( r \in (0,1) \) is common knowledge, while player 2's value \( \theta \in [0,1] \) is drawn from the uniform distribution and observed only by player 2.

\begin{enumerate}
    \item Is there any Bayesian Nash Equilibrium (BNE) in which player 1 always chooses to shirk? If there is, what is player 2's strategy in this BNE?
\begin{solution}
    Consider the \textbf{strategy in which player $1$ chooses to shirk always} and \textbf{player $2$ chooses to shirk regardless of type $\theta.$} Denote by $a_i$ the action taken by player $i.$ Note that $a_i \in \{W,S\}.$

    Fix player $2$ with this strategy. Then player $1$ receives a payoff from shirking of 
    \[U_1(a_1 = S | \theta)= U_1(a_1 = S) = 0.\] Player $1$ receives a payoff from working of 
    \[U_1(a_1 = W\mid \theta)  = U_1(a_1 = W)= -c.\] This will be optimal for player $1$ as long as $\boxed{0\geq -c},$ which is always true. 

\end{solution}
    \item Assume there is a BNE in which player 1 always chooses to work. What is player 2's strategy in this BNE?
\begin{solution}
 Fix player $1$ with this strategy. Then player $2$ receives a payoff from shirking of 
    \[U_2(a_2 = S \mid \theta) = U_2(a_2 = S) = 0.\] Player $2$ receives a payoff from working of 
    \[U_2(a_2 = W \mid \theta) = -c.\] This will be optimal for player $2$ as long as $0\geq -c,$ which is always true. 

    Thus, this strategy denotes a BNE. 
The strategy is for player $1$ to always work and for player $2$ to work if, and only if, $\theta \geq c$.

To see this, first fix player $1,$ then player 2's payoff matrix is given by \begin{table}[H]
        \centering
        \begin{tabular}{c|c|c}
             & Work & Shirk \\
             \hline
             Work& $\theta-c$& 0 \\
        \end{tabular}
    \end{table} 
the expected payoffs are given by 
\[\bbE[W_2] = \theta -c, \quad \bbE[S_2] = 0,\] and so player $2$ will work if, and only if, $\theta - c \geq 0 \implies \boxed{\theta \geq c.}$ That is, if the payoff of player $2$ is greater than the cost, then he will always work. 
    
\end{solution}
    \item Derive the condition in terms of \( c \) and \( r \) under which there exists a BNE in which player 1 always chooses to work.
\begin{solution}
Fix player $1$ $2$ with the strategies as above. Consider that since $\theta \sim U([0,1]),$ then if $c\in [0,1]$
\[\bbP\{\theta < c\} = c, \quad \bbP\{\theta \geq c\} = 1-c.\] Then the payoff player $1$ receives from working is
\[\bbE[W] = (r-c)\bbP\{\theta \geq c\} + (-c)\bbP\{\theta < c\} = (r-c)(1-c) -c^2 = r  - rc -c +c^2 - c^2 = r(1-c) - c.\] The payoff player $1$ receives from shirking is
\[\bbE[S_1]= 0.\] Thus, we need $r(1-c) \geq c$ and so player $1$ will work if $\boxed{r\geq \frac{c}{1-c}},$ which is always met, but by our strategy profile, player $2$ always works.

Thus, this strategy denotes a BNE. Consider when $r \geq \frac{c}{1-c}.$ Then the strategy profile where player $1$ always works and player $2$ works if, and only if, $\boxed{\theta \geq c}$ constitutes a BNE.
\end{solution}
    \item Assume there is a BNE in which player 1 chooses to work with probability \( p \in (0,1) \). What is player 2's strategy in this BNE?
\begin{solution}
    We claim that the strategy with which player $1$ chooses to work with probability $p$ and shirks with probability $1-p$ and player $2$ works if and only if $\theta \geq \frac{c}{p}$ constitutes a BNE. 

First, fix player $1$s strategy. Then the payoffs of player $2$ are given by 
\[\bbE[W_2] = (\theta - c)p + (-c)(1-p) = \theta p -cp - c + cp = \theta p -c\]
\[\bbE[S_2] = 0.\]
Thus, player $2$ will work if, and only if, 
\[\theta p -c \geq 0 \implies \boxed{\theta \geq \frac{c}{p}}.\] 


\end{solution}
    \item Derive the condition in terms of \( c \) and \( r \) under which there exists a BNE in which player 1 chooses to work with some probability \( p \in (0,1) \), and calculate \( p \) as a function of \( c \) and \( r \).
    
    (Hint: \( p \in (0,1) \) requires player 1 to be indifferent between working and shirking.)
\begin{solution}
Fix player $2$s strategy from above to work iff $\theta \geq \frac{c}{p}$ Then the payoffs of player $1$ are given by 
\[\bbE[W_1] = (r-c)\bbP\{\theta \geq \frac{c}{p}\} + (-c)\bbP\{\theta < \frac{c}{p}\} = (r-c)(1 - \frac{c}{p}) - \frac{c^2}{p} = r - \frac{rc}{p} - c\]
\[\bbE[S_1] = 0,\] and so player $1$ will work if, and only if, 
\[r - \frac{rc}{p} - c\geq 0 \iff rp - rc - cp \geq 0 \iff p(r-c) \geq rc \iff p \geq \frac{rc}{r-c}\] Thus, 
    \[\boxed{p = \frac{rc}{r-c}}\] is where player $1$ is indifferent between working and shirking. Since $p$ is a probability, we have that 
    \[0 < p < 1 \implies 0 < \frac{rc}{r-c} < 1 \implies rc < r-c \implies \boxed{\frac{c}{1-c} < r}\]
\end{solution}
\end{enumerate}

\newpage
\section*{Problem 2}
Consider an information structure with a prior $\mu \in \Delta(\Theta)$ and a type generating function $\pi: \Theta \to \Delta(\Theta).$ Prove that if this information structure is conditionally independent, then 
\[\pi(\theta, s_{-i} \mid s_i) = \pi(\theta \mid s_{i})\pi(s_{-i} \mid \theta) = \frac{\pi(s_{-i} \mid \theta) \pi(s_i \mid \theta)\mu(\theta)}{\sum_{\theta \in \Theta}\pi(s_i \mid \theta )\mu(\theta)}\]

\begin{solution}
For the second equality, consider that by Bayes' rule
\[\pi(\theta \mid s_i) =
\frac{\pi(\theta , s_i)}{\pi(s_i)}.\] For the numerator, we have that again, by Bayes' rule:
\[\pi(\theta, s_i) = \pi(s_i \mid \theta)\pi(\theta) = \pi(s_i \mid \theta)\mu(\theta).\] For the denominator, we use the law of total probability, 
\[\pi(s_i) = \sum_{\theta \in \Theta} \pi(s_i \mid \theta)\pi(\theta) = \sum_{\theta \in \Theta}\pi(s_i \mid \theta)\mu(\theta).\] Putting it all together, we see that 
\[\pi(\theta \mid s_i) =
\frac{\pi(s_i \mid \theta)\mu(\theta)}{\sum_{\theta \in \Theta}\pi(s_i \mid \theta)\mu(\theta)}.\] 

For the first equality, we again use Bayes' rule:
\begin{align*}
    \pi(\theta, s_{-i} \mid s_i) &= \frac{\pi(\theta, s_{-i}, s_i)}{\pi(s_i)}\\
    &= \frac{\pi(s_i)\frac{\pi(\theta, s_i)}{\pi(s_i)}\frac{\pi(\theta, s_{-i}, s_i)}{\pi(\theta, s_{i})}}{\pi(s_i)}\\
    &= \pi(\theta \mid s_i)\pi(s_{-i} \mid \theta, s_i)\\
    &= \pi(\theta \mid s_i) \frac{\pi(s_{-i}, \theta, s_i)}{\pi(\theta, s_i)}\\
    &= \pi(\theta \mid s_i) \frac{\pi(s_{-i} s_i \mid \theta)\pi(\theta)}{\pi(s_i \mid \theta)\pi(\theta)}\\
    &= \pi(\theta \mid s_i) \frac{\pi(s_{-i}  \mid \theta)\pi(s_i \mid \theta)\pi(\theta)}{\pi(s_i \mid \theta)\pi(\theta)}\\
    &= \pi(\theta \mid s_i)\pi(s_{-i} \mid \theta).
\end{align*}
Note that we used conditional indpendence for the second to last equality.
\end{solution}

\newpage
\section*{Problem 3}
\begin{problem}

Two armies, A and B, are about to fight against each other, each consists of a continuum of soldiers. Soldier in army A is indexed by \( i \in [0, 1] \) while soldier in army B is indexed by \( j \in [0, 1] \). Each soldier \( i \) chooses whether to fight for group A or to desert, \( a_i = 0,1 \), where \( a_i = 1 \) indicates fighting for group A. Each soldier \( j \) chooses whether to fight for group B or to desert, \( b_j = 0, 1 \), where \( b_j = 1 \) indicates fighting for group B. Fighting turnout in army A and B are respectively

\[
A = \int_0^1 a_i \, di, \quad B = \int_0^1 b_j \, dj
\]

A soldier's payoff depends on which army wins and on whether he/she fights for the army. Specifically, each \( i \) gets

\[
\begin{array}{c|c|c}
 & \text{A wins} & \text{B wins} \\
\hline
a_i = 1 & r_A - c_A & -c_A \\
a_i = 0 & 0 & 0 \\
\end{array}
\]

where \( c_A > 0 \) is the cost of fighting for army A and \( r_A > c_A \) is the reward. Note that a soldier in group A gets the reward iff he participates the fight and army A defeats army B. Similarly, each \( j \) gets

\[
\begin{array}{c|c|c}
 & \text{A wins} & \text{B wins} \\
\hline
b_j = 1 & -c_B & r_B - c_B \\
b_j = 0 & 0 & 0 \\
\end{array}
\]

where \( c_B > 0 \) and \( r_B > c_B \).

A defeats B iff

\[
A - B \geq \theta,
\]

where \( \theta \in \mathbb{R} \) measures the advantage of B against A due to factors other than soldier turnout.

Nobody knows \( \theta \) and the prior is uniform. Each \( i \) receives a signal

\[
s_i = \theta + \epsilon_i
\]

and each \( j \) receives a signal

\[
s_j = \theta + \epsilon_j,
\]

where \( \epsilon_i \) and \( \epsilon_j \) are i.i.d. and drawn from a c.d.f. \( F \).
\end{problem}

\begin{enumerate}
    \item (1 point) If \( \theta > 1 \), who would win and what would each soldier do? If \( \theta < -1 \), who would win and what would each soldier do?
\begin{solution}
    Suppose $\theta >1,$ then army $A$ will defeat $B$ iff $A - B \geq \theta.$ As defined, $\max A = 1$ and $\min B = 0,$ that is, the maximum $\max A - B = 1,$ and so $A - B < \theta$ for any possible combination of the armies. Thus, \textbf{army $A$ will lose no matter what.} 
    
    Given this, then for any $i \in [0,1],$ if $i$ fights, he receives a payoff of $-c_A$ with probability 1. If he surrenders, then he receives a payoff $0$ with probability $1.$ Since $0 \geq -c_{A},$ \textbf{then $i$ will surrender no matter the size of $B$.}

    For any $j \in [0,1],$ if $j$ fights, he receives a payoff of $r_B - c_B$ with probability $1.$ If he surrenders, then he receives a payoff of $0.$ Since $r_B - c_B > 0$ by assumption, \textbf{then $j$ will fight no matter what. }
\rule{\linewidth}{0.4pt}


Suppose $\theta <-1.$ Then since $\min A = 0$ and $\max B = 1,$ we have that $\min A - B = -1.$ Thus, $A-B \geq -1 >\theta$ no matter the combination of the armies, and so \textbf{$A$ will win no matter what.} Using symmetrical logic as the above case, \textbf{a soldier $i$ will fight unconditionally} and \textbf{a soldier $j$ will surrender unconditionally.}
\end{solution}
    \item (1 point) What is a critical state of this game?
\begin{solution}
    We may guess that army $A$ will win iff $\boxed{\theta  \leq \theta^*}$ for some $\theta^* \leq 1.$ This is because the lower $\theta$ is, the lower the advantage of $B$ against $A,$ and so $A - B$ has a higher chance of being above $\theta.$   
\end{solution}
    \item (2 points) Suppose we have a critical state \( \theta^* \). What would soldier \( i \) in army \( A \) do after observing \( s_i \)? What would soldier \( j \) in army \( B \) do after observing \( s_j \)?
\begin{solution}
    Fix army $B$ and consider soldier $i.$ He receives payoff of $r_A  - c_A$ when they win, and $-c_A$ when $A$ loses. The probability of $A$ winning (given the information that $i$ has) is (denoting a win by W and a loss by L), 
    \[\bbP\{W \mid s_i\} = \bbP\{\theta \leq \theta^* \mid s_i\} = \bbP\{s_i - \varepsilon_i \leq \theta^*\} = \bbP\{\varepsilon_i > s_i - \theta^*\} = 1- F(s_i - \theta^*) = 1-\bbP\{L\mid s_i\}.\] Thus, the payoff of $i$ going to war is 
    \[\bbE[a_i = 1 \mid s_i] = (r_A - c_A)(1- F(s_i - \theta^*)) + (-c_A)(F(s_i - \theta^*) = r_A(1- F(s_i - \theta^*) - c_A.\] The payoff of $i$ surrendering is
    \[\bbE[a_i = 0 \mid s_i] = 0.\] Thus, player $i$ fights in this war if and only if
    \begin{align*}
        &r_A(1- F(s_i - \theta^*) - c_A \geq 0\\
        &\implies F(s_i - \theta^*) \leq 1 - \frac{c_A}{r_A}\\
        &\implies s_i- \theta^* \leq F^{-1}(1 - \frac{c_A}{r_A})\\
        &\implies \boxed{s_i \leq \theta^* + F^{-1}(1 - \frac{c_A}{r_A})} =: k_A.
    \end{align*}
Thus, player $i$ fights iff $s_i \leq \theta^* + F^{-1}(1 - \frac{c_A}{r_A}) = k_A.$ Using symmetrical logic, player $j$ will fight iff 
\[\boxed{s_j \geq \theta^* + F^{-1}(\frac{c_B}{r_B})} =: k_B.\]
\end{solution}
    \item (2 points) Prove that given the strategies of \( i \)'s and \( j \)'s you derive in 2.3, there must be a critical state. (Hint: use the strong law of large numbers).
\begin{solution}
    We assume that each soldier makes their decision independently. The strong law of large numbers then states that 
    \[\int_0^1 a_i di = \bbE[a_i \mid \theta]\] Thus, using the definition of expectation,
    \begin{align*}
        A &= \int_0^a a_i di\\
        &= \bbE[a_i \mid \theta]\\
        &= \bbP\{a_i = 1 \mid \theta\}\\
        &= \bbP\{s_i \leq k_A \mid \theta\}\\
        &= \bbP\{\theta + \varepsilon_i \leq k_A \}\\
        &= F(k_A - \theta)
    \end{align*}
    Similarly, 
    \begin{align*}
        B &= \int_0^1 b_j dj\\
        &= \bbE[b_i \mid \theta]\\
        &= \bbP\{b_i = 1 \mid \theta\}\\
        &= \bbP\{s_j \geq k_B \mid \theta\}\\
        &= \bbP\{\theta + \epsilon_j \geq k_B\}
    \end{align*}
    Symmetrically, we have that
    \[B = 1 - F(k_B - \theta).\] Then army $A$ will win and $B$ will lose when 
    \[A - B \geq \theta \iff F(k_A - \theta) - (1 - F(k_B - \theta)) \geq \theta \iff  F(k_A - \theta) +F(k_B - \theta) \geq 1 + \theta.\]
    We claim there exists some $\theta^*$ such that 
    \[F(k_A - \theta^*) +F(k_B - \theta^*) = 1 + \theta^*.\] To see this, define an auxiliary function $\varphi: \Theta \to \bbR$ by
    \[\varphi(\theta):= F(k_A - \theta) +F(k_B - \theta) - 1 - \theta.\] It suffices to show that there is some $\theta^*$ such that $\varphi(\theta^*) = 0.$ Recall a few facts about any cdf:
    \begin{itemize}
        \item $F$ is monotonically increasing;
        \item $F$ is right continuous, but assume that it is continuous since we are dealing with continuous random variables;
        \item $\lim_{x\to -\infty}F(x) = 0,$ and $\lim_{x\to \infty} F(x) = 1.$ 
    \end{itemize}
    Thus, we have that $\varphi$ is continuous over $\Theta$ with 
    \[\lim_{\theta \to -\infty} \varphi(\theta)= \lim_{\theta \to -\infty} F(k_A - \theta) + \lim_{\theta \to -\infty} F(k_B - \theta) - 1 - \lim_{\theta \to -\infty} \theta = 1 + 1 - 1 + \infty = \infty.\] Similarly, 
    \[\lim_{\theta \to \infty} \varphi(\theta)=-\infty\]
    By the IVT, there is some $\theta^* \in \Theta$ such that $\varphi(\theta^*) = 0.$ That is,
    \[F(k_A - \theta^*) +F(k_B - \theta^*) = 1 + \theta^*.\]
    
\end{solution}
    \item (2 points) Solve the critical state in BNE and discuss how it relates to \( r_A, c_A \) and \( r_B, c_B \). Comment on the solution as if you were commanding an army.
\begin{solution}
    We have 3 equations, 3 unknowns:
    \begin{align}
        F(k_A - \theta^*) +F(k_B - \theta^*) = 1 + \theta^*
    \end{align}
    \begin{align}
        \theta^* + F^{-1}(1 - \frac{c_A}{r_A}) = k_A
    \end{align}
    \begin{align}
        \theta^* + F^{-1}(\frac{c_B}{r_B}) = k_B
    \end{align}
    Plug in (2) and (3) into (1):
    \begin{align*}
        F(k_A - \theta^*) +F(k_B - \theta^*) &= F(\theta^* + F^{-1}(1 -\frac{c_A}{r_A}) - \theta^*) +F( \theta^* + F^{-1}(\frac{c_B}{r_B}) - \theta^*)\\
        &= F(F^{-1}(1 - \frac{c_A}{r_A})) + F(F^{-1}(\frac{c_B}{r_B}))\\
        &= 1 - \frac{c_A}{r_A} + \frac{c_B}{r_B}\\
        &= 1 + \theta^*
    \end{align*}
    Thus,
    \[\boxed{\theta^* = \frac{c_B}{r_B}- \frac{c_A}{r_A}} \]

Suppose I am commander A. I would want $\theta^*$ to be as high as possible, so then I would want $\frac{c_A}{r_A}$ to be as low as possible. Thus, I would start a cult that glorifies death in war (much like the culture surrounding war in WW2 Japan) to decrease the cost of fighting, and I would pay my soldiers more! Commander $B$ would do the same with his soldiers, since he wants $\theta^*$ to be low.
\end{solution}
\end{enumerate}



\end{document}